\documentclass[11pt]{article}

\usepackage[a4paper,margin=1in]{geometry}
\usepackage[T1]{fontenc}
\usepackage[utf8]{inputenc}
\usepackage{lmodern}
\usepackage{microtype}
\usepackage{graphicx}
\usepackage{booktabs}
\usepackage{array}
\usepackage{longtable}
\usepackage{hyperref}
\usepackage{xcolor}
\usepackage{enumitem}
\usepackage{float}

\hypersetup{
    colorlinks=true,
    linkcolor=blue,
    urlcolor=blue,
    citecolor=blue
}

\setlength{\parskip}{0.5em}
\setlength{\parindent}{0pt}

\newcommand{\courseworktitle}{Patronising and Condescending Language Detection}
\newcommand{\repositoryurl}{https://github.com/J1mmyQi/pcl-detection-nlp}
\newcommand{\studentname}{Jinming Qi}

\newcommand{\optionalfigure}[4]{
    \begin{figure}[H]
        \centering
        \IfFileExists{#1}{
            \includegraphics[width=#2]{#1}
        }{
            \fbox{
                \parbox[c][2.2in][c]{0.85\linewidth}{
                    \centering
                    Placeholder for figure:\\
                    \texttt{#1}
                }
            }
        }
        \caption{#3}
        \label{#4}
    \end{figure}
}

\begin{document}

\begin{titlepage}
    \centering
    {\Large \textbf{\courseworktitle} \par}
    \vspace{1.2em}
    {\large A coursework report in research-paper format \par}
    \vspace{2em}

    \begin{tabular}{>{\raggedright\arraybackslash}p{0.28\linewidth} p{0.62\linewidth}}
        Author & \studentname \\
        Repository & \href{\repositoryurl}{\repositoryurl} \\
        Task & SemEval 2022 Task 4, Subtask 1 (binary PCL classification) \\
    \end{tabular}

    \vfill

    {\small
    This report is structured to match the coursework specification while using a
    concise research-paper presentation style. The content below currently
    reflects the currently retained experiment set and the latest local
    results available in the repository.
    \par}
\end{titlepage}


    \section{Introduction}
    Patronising and condescending language (PCL) detection is difficult because the
    relevant signal lies in framing rather than in topic alone. Paragraphs about
    refugees, migrants, women, or poor families may be neutral, supportive, or
    condescending while using broadly similar vocabulary, so the task depends on
    distinguishing rhetorical stance rather than simply identifying subject matter.

    This report presents a compact research workflow for the coursework task: it
    reviews the task literature, analyzes the dataset, motivates the modeling
    approach, and then evaluates the strongest currently retained local runs on the
    official development split. The repository now contains tuned sparse baselines,
    RoBERTa-based runs, and score-level combinations of saved models. The emphasis
    here is on documenting what these local experiments show, rather than claiming
    that one architecture is already the final endpoint of the coursework.


    \section{Critical Review of the Task Paper}

    \subsection{Research problem and motivation}
    The reviewed work asks whether NLP systems can detect patronising and
    condescending language in paragraphs about vulnerable communities. This is
    challenging because the distinction is primarily pragmatic: the same topics may
    appear in both neutral reporting and PCL, while the label depends on framing,
    tone, and implied agency. The task is therefore a meaningful benchmark for
    models that claim to capture stance rather than only topic.

    \subsection{Main contributions of the reviewed paper}
    The task paper and the earlier dataset paper make four practical contributions
    \cite{semeval2022,baseline}: they define a shared supervised task, release a
    curated dataset, provide both binary and fine-grained label settings, and
    establish baseline systems with a common evaluation protocol. Their most useful
    contribution for this project is the operationalization of a socially subtle
    phenomenon into a reproducible benchmark with fixed splits, submission files,
    and a clear primary metric.

    \subsection{What the paper does well}
    Its main strength is benchmark design. The task is difficult for the right
    reason: PCL is often subtle, context-dependent, and superficially similar to
    sympathetic language. The train/dev/test split supports fair comparison, while
    the binary setting keeps the baseline problem tractable without making it easy.
    The connection to a richer label taxonomy also gives the benchmark value beyond
    the shared-task setting.

    \subsection{Limitations and open issues}
    The key limitation is annotation subjectivity: the boundary between empathy and
    condescension is often ambiguous, so some label noise is unavoidable. A second
    limitation is that strong scores can still come from shallow lexical shortcuts,
    especially when certain community keywords recur. Finally, the benchmark
    measures a controlled proxy for the problem rather than the full social
    phenomenon, so leaderboard performance should not be treated as complete
    evidence of discourse-level understanding.

    \subsection{How the literature review informed this project}
    The review directly informed three decisions in this project: selecting
    positive-class F1 as the main metric, using class-balanced training to address
    imbalance, and extending a word-only TF-IDF baseline with character n-grams to
    capture some stylistic variation. It also motivates keeping contextual modeling
    in scope once the gains from purely sparse features begin to saturate.


    \section{Dataset Analysis and Preprocessing}

    \subsection{Dataset overview}
    The coursework uses the \texttt{dontpatronizeme\_pcl.tsv} data together with the
    official train, development, and test splits. The implemented loader converts
    the original ordinal labels into the binary task required by the specification:
    labels \texttt{0} and \texttt{1} are mapped to \texttt{No PCL}, while labels
    \texttt{2}, \texttt{3}, and \texttt{4} are mapped to \texttt{PCL}.

    Table~\ref{tab:data-summary} summarizes the currently loaded splits.

    \begin{table}[H]
    \centering
    \caption{Summary of the available splits after binary conversion}
    \label{tab:data-summary}
    \begin{tabular}{lrrrrr}
    \toprule
    Split & Size & Positives & Negatives & Avg. tokens & Max tokens \\
    \midrule
    Train & 8375 & 794 & 7581 & 48.68 & 909 \\
    Dev   & 2094 & 199 & 1895 & 47.41 & 272 \\
    Test  & 3832 & -- & -- & 48.36 & 379 \\
    \bottomrule
    \end{tabular}
    \end{table}

    The most important immediate observation is the severe class imbalance. The
    positive class accounts for roughly 9.5\% of both the train and development
    splits. This means accuracy would be a misleading primary metric and justifies
    the coursework's emphasis on the positive-class F1 score.

    \subsection{EDA technique 1: class distribution analysis}
    The first EDA step is class-distribution profiling. Because the dataset is
    strongly skewed toward the negative class, the current classifier uses
    \texttt{class\_weight="balanced"} so minority-class errors matter more during
    training.

    \optionalfigure
    {report_assets/eda_label_distribution.png}
    {0.7\linewidth}
    {Binary label distribution in the training split.}
    {fig:label-distribution}

    \subsection{EDA technique 2: keyword frequency analysis}
    The second EDA technique is lexical frequency analysis over topic keywords. In
    the positive training subset, the most frequent keywords are
    \texttt{homeless}, \texttt{in-need}, \texttt{poor-families}, and
    \texttt{hopeless}. These keywords are informative but not decisive: they mark
    topic domains, while the label still depends on how those groups are framed.

    \begin{table}[H]
    \centering
    \caption{Top keywords in positive training examples}
    \label{tab:positive-keywords}
    \begin{tabular}{lr}
    \toprule
    Keyword & Count \\
    \midrule
    homeless & 149 \\
    in-need & 143 \\
    poor-families & 112 \\
    hopeless & 98 \\
    refugee & 73 \\
    disabled & 67 \\
    vulnerable & 60 \\
    women & 38 \\
    migrant & 31 \\
    immigrant & 23 \\
    \bottomrule
    \end{tabular}
    \end{table}

    \optionalfigure
    {report_assets/top_positive_keywords.png}
    {0.8\linewidth}
    {Most frequent topic keywords in positive training examples.}
    {fig:keyword-frequency}

    This observation motivates a mixed feature space: lexical features are retained
    because they carry signal, but character n-grams are added to capture recurring
    local phrasing patterns beyond simple topic words.

    \subsection{Preprocessing decisions}
    The preprocessing pipeline is intentionally light:
    \begin{itemize}[leftmargin=1.5em]
    \item labels are converted to the binary task required by the specification;
    \item Unicode accents are normalized during TF-IDF feature extraction;
    \item stop-word removal is disabled in the current best configuration;
    \item word-level and character-level TF-IDF features are both enabled;
    \item no heavy text normalization is applied, so most stylistic cues remain available.
    \end{itemize}

    This is deliberate. Heavy normalization could erase the same stylistic cues that
    help distinguish neutral reporting from soft condescension.


    \section{Proposed Approach}

    \subsection{Method summary}
    The repository now explores three related model families rather than a single
    fixed classifier:
    \begin{itemize}[leftmargin=1.5em]
    \item tuned sparse models built from word and character TF-IDF features;
    \item RoBERTa-based checkpoints, including a fine-tuned model and a lighter probe-style variant;
    \item score-level fusion over saved artifacts so lexical and contextual signals can be compared and combined.
    \end{itemize}

    \subsection{Rationale}
    This progression addresses a specific weakness of the plain baseline. Word
    n-grams capture topic and short lexical cues, but many PCL examples differ
    mainly in framing. Sparse models therefore provide a transparent baseline,
    while contextual encoders and score fusion test whether broader sentence-level
    representations and complementary decision boundaries recover additional
    positive cases.

    \subsection{Expected outcome}
    The target was to improve recall on PCL examples without collapsing precision.
    That objective was achieved locally: the stronger runs recover substantially
    more positives on the development split than the original sparse baseline. At
    the same time, this report treats the current top local run as evidence from an
    ongoing comparison, not as a claim that architectural exploration is finished.


    \section{Implementation and Training Setup}

    \subsection{System description}
    The codebase is organized into three functional layers:
    \begin{itemize}[leftmargin=1.5em]
    \item \texttt{data\_pipeline.py}: loading, preprocessing, and EDA;
    \item \texttt{models.py}: model definitions;
    \item \texttt{training\_pipeline.py}: training, evaluation, and export.
    \end{itemize}

The companion notebook (\texttt{coursework\_runner.ipynb}) is the main
interactive runner, while the Python modules remain the implementation source
of truth. Historical exploratory runs are retained separately under
\texttt{archive/} so that the active repository surface stays focused on the
final deliverable path. The \texttt{BestModel/} bundle now also includes a
direct notebook copy, a code snapshot, mirrored submission files, and the
checkpoint files needed to reload the selected final ensemble without relying
on the gitignored \texttt{artifacts/} directory.

    \subsection{Best configuration}
    Table~\ref{tab:best-config} summarizes the current best configuration stored in
    \texttt{configs/best\_artifact\_ensemble.json}.

    \begin{table}[H]
    \centering
    \caption{Current highest-scoring local run configuration}
    \label{tab:best-config}
    \begin{tabular}{ll}
    \toprule
    Component & Setting \\
    \midrule
    Member 1 & Fine-tuned RoBERTa artifact from \texttt{artifacts/roberta\_download} \\
    Member 2 & RoBERTa probe artifact from \texttt{artifacts/best\_roberta\_probe} \\
    Member 3 & Best non-transformer ensemble from \texttt{artifacts/best\_hybrid\_ensemble} \\
    Member 4 & Fine-tuned probe variant from \texttt{artifacts/roberta\_probe\_finetuned\_01} \\
    Score normalization & Per-member z-score normalization before fusion \\
    Ensemble weights & 0.30 / 0.38 / 0.12 / 0.20 \\
    Decision rule & weighted score average with threshold \texttt{1.1838} \\
    \bottomrule
    \end{tabular}
    \end{table}

    \subsection{Development-set results}
    The current best development-set metrics are shown in
    Table~\ref{tab:dev-results}. These values were read from the latest
    \texttt{artifacts/best\_artifact\_ensemble/metrics.json} file.

    \begin{table}[H]
    \centering
    \caption{Development-set performance of the current best model}
    \label{tab:dev-results}
    \begin{tabular}{lrrrrrrr}
    \toprule
    Precision & Recall & F1 & TP & FP & FN & TN \\
    \midrule
    0.6009 & 0.6432 & 0.6214 & 128 & 85 & 71 & 1810 \\
    \bottomrule
    \end{tabular}
    \end{table}

    \begin{table}[H]
    \centering
    \caption{Comparison across representative local runs}
    \label{tab:experiment-comparison}
    \begin{tabular}{lrrr}
    \toprule
    Configuration & Precision & Recall & F1 \\
    \midrule
    Initial baseline (\texttt{tfidf\_svm}) & 0.4559 & 0.3116 & 0.3701 \\
    Best non-transformer ensemble & 0.4075 & 0.5427 & 0.4655 \\
    Strongest single fine-tuned RoBERTa & 0.6236 & 0.5578 & 0.5889 \\
    Current highest-scoring local run & 0.6009 & 0.6432 & 0.6214 \\
    \bottomrule
    \end{tabular}
    \end{table}

    Relative to the original local baseline, the current highest-scoring run
    improves the development-set F1 score by roughly 0.251 absolute points. The
    final fused run remains recall-oriented, but the latest calibration also
    recovers some precision relative to the earlier ensemble version, reducing the
    false-positive burden while still identifying substantially more positive PCL
    cases than the baseline. For this task, that trade-off is reasonable because
    the positive class is rare and under-detection is a major failure mode.

The comparison in Table~\ref{tab:experiment-comparison} also serves as a small
ablation-style local evaluation: it shows first that sparse tuning materially
improves on the baseline, then that a full contextual checkpoint outperforms
the non-transformer family, and finally that score-level fusion of
complementary saved runs adds another incremental gain.

\optionalfigure
    {report_assets/dev_precision_recall_curve.png}
    {0.6\linewidth}
    {Precision-recall curve for the current best development-set run.}
    {fig:pr-curve}

Figure~\ref{fig:pr-curve} complements the single-threshold F1 score by showing
the full precision-recall trade-off over the development split. The area under
this curve is also strong in local terms (average precision 0.6091). The shape
confirms that the retained operating point is part of a broader high-recall
region rather than an isolated threshold artifact, which makes the local result
more credible as an evaluation finding.

\optionalfigure
    {report_assets/dev_confusion_matrix.png}
    {0.55\linewidth}
    {Confusion matrix for the best development-set run.}
    {fig:confusion-matrix}

    The current model is stronger than the earlier untuned baseline used at the
    start of the project and now exceeds the coursework's local development
    reference threshold. Its main value in this report is therefore as evidence
    that combining contextual and lexical signals materially improves local
    performance, while still leaving room for stronger validation and further
    checkpoint diversity.

    \subsection{Why the tuned model performs better}
    The improvement comes from combining complementary score sources rather than
    relying on a single decision surface. The fine-tuned RoBERTa artifact captures
    broader contextual framing, the two probe-style variants add lighter but
    diverse contextual signals, and the non-transformer ensemble preserves useful
    lexical bias toward topic and local phrasing. Together they produce a less
    brittle recall profile than any one member alone.

    \subsection{Error analysis}
    The development-set errors follow a clear pattern: many false negatives are
    supportive or policy-oriented references to vulnerable groups that lack an
    overt lexical cue. Representative failure modes include:
    \begin{itemize}[leftmargin=1.5em]
    \item factual refugee and migrant reporting that remains close to the PCL decision boundary;
    \item long news-style sentences where the key framing cue is diluted by surrounding context;
    \item examples about women or poor families where condescension depends on tone rather than explicit wording.
    \end{itemize}

    This suggests that the current sparse linear model still underfits
    ensemble still faces a discourse-level boundary problem: even stronger runs can
    identify topic correctly before missing the final rhetorical distinction.

    \begin{table}[H]
    \centering
    \caption{Representative qualitative errors from the development set}
    \label{tab:error-examples}
    \begin{tabular}{p{0.12\linewidth} p{0.16\linewidth} p{0.62\linewidth}}
    \toprule
    ID & Keyword & Observation \\
    \midrule
    8330 & refugee & A short sentence about refugee resettlement is predicted as non-PCL, showing that neutral policy language near the decision boundary is still often missed. \\
    1572 & migrant & A sentence about dignity and disdain toward migrant workers is predicted as non-PCL, suggesting the system under-detects more implicit evaluative framing. \\
    5506 & vulnerable & A sentence about helping ``vulnerable families'' is predicted as non-PCL, indicating that positive institutional language can still carry soft framing cues the current local run does not reliably capture. \\
    \bottomrule
    \end{tabular}
    \end{table}

    False positives show the complementary failure mode. Once recall is increased,
    the model becomes more willing to label emotionally charged language as PCL even
    when the framing is not clearly condescending. This remains the main trade-off
    of the current setup: broader positive coverage makes the decision boundary
    harder to maintain even after combining multiple complementary models.

    \subsection{Implementation notes}
    The repository now validates submission formatting before writing prediction
    files. This ensures that:
    \begin{itemize}[leftmargin=1.5em]
    \item each file contains exactly one prediction per input example;
    \item each line contains only \texttt{0} or \texttt{1};
    \item the exported files are directly compatible with the coursework
    submission requirements.
    \end{itemize}


    \section{Conclusion}
    This report documents a complete local PCL detection pipeline in a format that
    matches the coursework requirements while retaining a compact research-paper
    presentation. The repository now supports a meaningful progression from sparse
    baselines to contextual runs and score-level fusion, and the strongest current
    local result improves clearly over the earlier baseline on the development
    split.

    The experimental gains are real, but the more useful conclusion at this stage
    is methodological caution: the local comparison now favors combining
    complementary signals, yet the remaining errors still show how difficult it is
    to separate empathy from condescension in borderline cases. The present report
    therefore treats the top local result as strong evidence, not as a final claim
    that model exploration is complete.

    \subsection*{Next technical step}
    The most justified next step is to make the validation protocol stronger rather
    than merely broader: additional full fine-tuning runs with different random
    seeds, cleaner calibration splits, and more disciplined checkpoint ensembling
    would provide a firmer estimate of what genuinely transfers beyond the current
    development set.


\section*{Submission notes}
\begin{itemize}[leftmargin=1.5em]
    \item Confirm that the title-page author and repository details match the final submitted repository state.
    \item The notebook saves report figures directly into \texttt{report\_assets/}; keep those committed copies so the compiled report includes the final plots.
    \item Preserve \texttt{dev.txt} and \texttt{test.txt} in their submission-ready form.
\end{itemize}


    \section*{References}
    \begin{thebibliography}{9}

    \bibitem{semeval2022}
    SemEval 2022 Task 4: Patronizing and Condescending Language Detection.\\
    Task website: \url{https://sites.google.com/view/pcl-detection-semeval2022/}

    \bibitem{baseline}
    Perez-Almendros, C. et al. ``You Don't Look Poor'': A Dataset and Baselines for
    Contextualized Detection of Patronizing and Condescending Language.\\
    ACL Anthology entry: \url{https://aclanthology.org/2020.coling-main.518/}

    \end{thebibliography}

    \end{document}
